% ============================================================================
% Section 4.2 — Simplifying Expressions by Combining Like Terms
% CCSS 7.EE.A.1
% ============================================================================
\section{Combining Like Terms}

\begin{stepGoal}
\begin{itemize}
    \item Identify like terms in an expression
    \item Simplify expressions by combining like terms
\end{itemize}
\end{stepGoal}

\begin{stepVocabBox}
    \vocabItem{Like Terms}{Terms with the \textbf{same variable} raised to the same power (e.g.\ $3x$ and $7x$).}
    \vocabItem{Coefficient}{The number in front of a variable. In $7x$, the coefficient is $7$.}
\end{stepVocabBox}

\begin{stepsCard}[How to Combine Like Terms]
    \stepItem{\textbf{Identify} like terms --- same variable, same exponent. Constants are like terms with each other.}
    \stepItem{\textbf{Group} like terms together (rewrite if needed).}
    \stepItem{\textbf{Add or subtract} the coefficients of each group. Keep the variable the same.}
\end{stepsCard}

% --- Worked Example 1 ---
\begin{stepExample}{Simplify $4x + 3 - 2x + 8$.}
    \stepShow{1} Like terms: $4x$ and $-2x$ (both have $x$); $3$ and $8$ (both constants).

    \stepShow{2} Group: $(4x - 2x) + (3 + 8)$.

    \stepShow{3} Combine coefficients: $4 - 2 = 2$, so $2x$. Constants: $3 + 8 = 11$.
    \stepResult{$2x + 11$}
\end{stepExample}

% --- Worked Example 2 ---
\begin{stepExample}{Simplify $\frac{1}{2}n + \frac{3}{4}n - 2$.}
    \stepShow{1} Like terms: $\frac{1}{2}n$ and $\frac{3}{4}n$ (both have $n$). The $-2$ has no like term.

    \stepShow{2} Group: $\left(\frac{1}{2}n + \frac{3}{4}n\right) - 2$.

    \stepShow{3} Common denominator: $\frac{2}{4}n + \frac{3}{4}n = \frac{5}{4}n$.
    \stepResult{$\frac{5}{4}n - 2$}
\end{stepExample}

\watchOut{$3x + 5$ cannot be simplified further --- $3x$ is a variable term and $5$ is a constant. They are NOT like terms!}

% ============================================================================
% PRACTICE
% ============================================================================
\newpage
\begin{stepPractice}{Combining Like Terms Practice}
    \resetProblems

    \stepPracticeHeader[step1Color]{Identify Like Terms}
    \prob Circle the like terms in $6a + 3b - 2a + 7$. \answerBlank[3cm]
    \answer{$6a$ and $-2a$ are like terms}

    \stepPracticeHeader[step2Color]{Group and Combine}
    \begin{multicols}{2}
    \prob Simplify $6a + 2a - 3$. \answerBlank[2.5cm]
    \answer{$8a - 3$}
    \prob Simplify $5x - 9 + 3x + 4$. \answerBlank[2.5cm]
    \answer{$8x - 5$}
    \end{multicols}

    \stepPracticeHeader[step3Color]{Combine with Negatives and Fractions}
    \prob Simplify $-2y + 7 + 6y - 10$. \answerBlank[2.5cm]
    \answer{$4y - 3$}

    \prob Simplify $\frac{2}{3}m - 5 + \frac{1}{3}m + 2$. \answerBlank[2.5cm]
    \answer{$m - 3$}

    \wordProblem{Sam has $3x + 8$ baseball cards and gives away $x + 2$ cards. Write a simplified expression for how many cards Sam has now.}{cards}
    \answer{$2x + 6$}
\end{stepPractice}
